\section{Discussion}

%Presents a clear statement and insightful interpretation of the results, accurately relates results to the hypothesis, and discusses any limitations or sources of error in the experiment.
%Questions are answered completely and correctly.

By examining the magnetic response of a steel toroid, we produced a plot of a hysteresis curve.
This meant that the direction of the applied current had a direct effect on the behaviour in which the material was
magnetised, and that magnetisation of steel can be intentionally adjusted to meet specifications.
The full magnetic response was a closed sigmoid-shaped loop, where calculating the area produced a value of
\qty{439.8}{\tesla\ampere\per\meter}, which described the work done per unit volume by the driving current for the
full cycle.
By applying thermodynamic equations, we determined that one cycle should increase the temperature of the
steel by \qty{1.34E-4}{\celsius}, but this would not yield a noticeable difference in temperature if
measured by a $\qty{0.1}{\celsius}$ resolution thermometer.
Additionally, this is the quantity only if we assume all the work done was converted directly
into
Joule heating, however this was not the case.
The experiment does not consider the thermal conduction and heat dissipation of the materials.

The maximum slope of the initial phase was used to derive the maximum magnetic susceptibility, $\qty{82.40(0.45)}{}$, a unitless quantity which describes how well the material aligns with the magnetic field.

The magnetic field inside an ideal toroid is constant, but in this experiment, we measured the field effect at a
cut-out portion.
This lead to the measured field intensity being slightly less than what it would normally be, had the toroid been an
ideal shape and complete.

Since there was some residual magnetisation left over from the material being previously magnetised, there was some
influence on the measured field strength.
Had the steel been completely demagnetised, the hysteresis curve would have better reflected the natural response of
the
material from an unbiased direction.


\section{Conclusion}

%Presents a clear, concise, focused, insightful and correct connection between the results and the investigated concepts.
%Explains why the findings are important, what they show, how they contribute to the broader understanding of the topic and recommendations for future research or improvements to
%experimental design.
%The section brings closure to the report.

This study demonstrated the magnetic response of a steel toroid depends on the applied current direction, and the
behaviour is different when ramping it up or down.
By integrating the shape the hysteresis curve produces, the work done was derived and used to calculate the increase
in temperature of the material.
The magnetic susceptibility was found to be a unique unitless quantity of the steel, which described how easily the
material was magnetised when the magnetic field was applied.
Some residual magnetisation from previous current cycles affected the magnetic response, but nonetheless the
experiment showed the expected properties of ferromagnetic materials.
